% ============================================================
% Chapitre 9 : Regression Multiple et ANOVA
% ============================================================
\chapter{R\'egression Lin\'eaire Multiple et ANOVA}
\label{ch:regression-multiple}

\begin{center}
\textit{<<~En science, la r\'ealit\'e est rarement unidimensionnelle.~>>}
\end{center}

% ------------------------------------------------------------
\section{Exemple motivant}
% ------------------------------------------------------------

On veut pr\'edire le prix d'un appartement en fonction de la surface, du nombre de pi\`eces, de l'\'etage et de la distance au centre-ville. Un seul pr\'edicteur ne suffit plus : il faut la r\'egression \textbf{multiple}.

% ------------------------------------------------------------
\section{Le mod\`ele matriciel}
% ------------------------------------------------------------

\begin{definition}[Mod\`ele lin\'eaire g\'en\'eral]
\index{regression lineaire@r\'egression lin\'eaire!multiple}
\[
\mathbf{Y} = \mathbf{X}\bm{\beta} + \bm{\varepsilon},
\]
o\`u $\mathbf{Y}\in\R^n$, $\mathbf{X}\in\R^{n\times p}$ (matrice de design, rang $p$), $\bm{\beta}\in\R^p$, $\bm{\varepsilon}\sim\mathcal{N}(\mathbf{0},\sigma^2\mathbf{I}_n)$.

Explicitement : $Y_i=\beta_0+\beta_1 x_{i1}+\cdots+\beta_{p-1}x_{i,p-1}+\varepsilon_i$.
\end{definition}

% ------------------------------------------------------------
\section{Estimation MCO}
% ------------------------------------------------------------

\begin{theoreme}[Solution MCO]
L'estimateur MCO est :
\[
\hat{\bm{\beta}} = (\mathbf{X}^\top\mathbf{X})^{-1}\mathbf{X}^\top\mathbf{Y}.
\]
Il minimise $\|\mathbf{Y}-\mathbf{X}\bm{\beta}\|^2$.
\end{theoreme}

\begin{proof}
$Q(\bm{\beta})=(\mathbf{Y}-\mathbf{X}\bm{\beta})^\top(\mathbf{Y}-\mathbf{X}\bm{\beta})$. La condition $\nabla_{\bm{\beta}}Q=\mathbf{0}$ donne $-2\mathbf{X}^\top(\mathbf{Y}-\mathbf{X}\bm{\beta})=\mathbf{0}$, soit $\mathbf{X}^\top\mathbf{X}\bm{\beta}=\mathbf{X}^\top\mathbf{Y}$ (\textbf{\'equations normales}).
\end{proof}

\begin{theoreme}[Propri\'et\'es]
\begin{enumerate}
\item $\E[\hat{\bm{\beta}}]=\bm{\beta}$ (sans biais).
\item $\Cov(\hat{\bm{\beta}})=\sigma^2(\mathbf{X}^\top\mathbf{X})^{-1}$.
\item $\hat{\sigma}^2=\frac{\|\mathbf{Y}-\mathbf{X}\hat{\bm{\beta}}\|^2}{n-p}$ est sans biais.
\item $\hat{\bm{\beta}}\sim\mathcal{N}(\bm{\beta},\sigma^2(\mathbf{X}^\top\mathbf{X})^{-1})$.
\end{enumerate}
\end{theoreme}

\begin{definition}[Matrice de projection]
\index{matrice de projection}
La \textbf{matrice chapeau} est $\mathbf{H}=\mathbf{X}(\mathbf{X}^\top\mathbf{X})^{-1}\mathbf{X}^\top$. On a $\hat{\mathbf{Y}}=\mathbf{H}\mathbf{Y}$ et $\mathbf{e}=(\mathbf{I}-\mathbf{H})\mathbf{Y}$.
\end{definition}

% ------------------------------------------------------------
\section{Inf\'erence}
% ------------------------------------------------------------

\subsection{Test individuel}

Pour $H_0:\beta_j=0$ :
\[
T_j = \frac{\hat{\beta}_j}{\mathrm{SE}(\hat{\beta}_j)} \sim t(n-p), \quad \mathrm{SE}(\hat{\beta}_j)=\hat{\sigma}\sqrt{[(\mathbf{X}^\top\mathbf{X})^{-1}]_{jj}}.
\]

\subsection{Test global de Fisher}

$H_0:\beta_1=\cdots=\beta_{p-1}=0$ :
\[
F = \frac{(\mathrm{SCT}-\mathrm{SCR})/(p-1)}{\mathrm{SCR}/(n-p)} = \frac{R^2/(p-1)}{(1-R^2)/(n-p)} \sim F(p-1,n-p).
\]

\begin{definition}[$R^2$ ajust\'e]
\index{R2 ajuste@$R^2$ ajust\'e}
\[
R^2_{\mathrm{adj}} = 1-\frac{n-1}{n-p}(1-R^2).
\]
P\'enalise l'ajout de variables inutiles.
\end{definition}

\subsection{Test de sous-mod\`ele}

Pour comparer un mod\`ele complet ($p$ param\`etres) et un mod\`ele r\'eduit ($q<p$ param\`etres) :
\[
F = \frac{(\mathrm{SCR}_{\text{r\'ed}}-\mathrm{SCR}_{\text{comp}})/(p-q)}{\mathrm{SCR}_{\text{comp}}/(n-p)} \sim F(p-q,n-p).
\]

% ------------------------------------------------------------
\section{ANOVA \`a un facteur}
% ------------------------------------------------------------

\begin{definition}[Mod\`ele ANOVA]
\index{ANOVA}
Soit $k$ groupes avec $n_j$ observations dans le groupe $j$ :
\[
Y_{ij} = \mu + \alpha_j + \varepsilon_{ij}, \quad i=1,\ldots,n_j,\; j=1,\ldots,k,
\]
avec $\sum n_j\alpha_j=0$ et $\varepsilon_{ij}\iid\mathcal{N}(0,\sigma^2)$.

$H_0:\alpha_1=\cdots=\alpha_k=0$ (toutes les moyennes sont \'egales).
\end{definition}

\begin{theoreme}[Tableau ANOVA]
\begin{center}
\begin{tabular}{lcccl}
\toprule
\textbf{Source} & \textbf{SC} & \textbf{ddl} & \textbf{CM} & \textbf{$F$} \\
\midrule
Facteur & $\mathrm{SCE}=\sum n_j(\bar{Y}_j-\bar{Y})^2$ & $k-1$ & $\mathrm{CME}$ & $F=\frac{\mathrm{CME}}{\mathrm{CMR}}$ \\
R\'esidu & $\mathrm{SCR}=\sum\sum(Y_{ij}-\bar{Y}_j)^2$ & $n-k$ & $\mathrm{CMR}$ & \\
Total & $\mathrm{SCT}=\sum\sum(Y_{ij}-\bar{Y})^2$ & $n-1$ & & \\
\bottomrule
\end{tabular}
\end{center}
Sous $H_0$ : $F\sim F(k-1,n-k)$.
\end{theoreme}

\begin{intuition}
L'ANOVA est un cas particulier de r\'egression multiple avec des variables explicatives indicatrices (dummy variables). Le test $F$ de l'ANOVA est exactement le test global de Fisher du mod\`ele correspondant.
\end{intuition}

\subsection{Comparaisons multiples}

\begin{attention}
Si $H_0$ est rejet\'ee, on veut savoir \emph{quels} groupes diff\`erent. Les tests de comparaison deux \`a deux sans correction augmentent le risque d'erreur de type~I global. On utilise :
\begin{itemize}
\item \textbf{Bonferroni} : $\alpha_{\mathrm{corrig\'e}}=\alpha/m$ o\`u $m=\binom{k}{2}$.
\item \textbf{Tukey HSD} : bas\'e sur la loi de l'\'etendue studentis\'ee.
\end{itemize}
\end{attention}

% ------------------------------------------------------------
\section{Impl\'ementation Python}
% ------------------------------------------------------------

\begin{python}
\begin{minted}{python}
import numpy as np
import pandas as pd
import statsmodels.api as sm
from statsmodels.formula.api import ols
from scipy import stats
import matplotlib.pyplot as plt

# --- Regression multiple ---
np.random.seed(42)
n = 50
surface = np.random.uniform(25, 120, n)
pieces = np.random.randint(1, 6, n)
etage = np.random.randint(0, 10, n)
prix = 30 + 2.5*surface + 15*pieces + 3*etage + np.random.normal(0, 15, n)

df = pd.DataFrame({'prix': prix, 'surface': surface,
                    'pieces': pieces, 'etage': etage})

# Ajustement
X = sm.add_constant(df[['surface', 'pieces', 'etage']])
model = sm.OLS(df['prix'], X).fit()
print(model.summary())

# R^2 ajuste
print(f"\nR^2 = {model.rsquared:.4f}")
print(f"R^2 ajuste = {model.rsquared_adj:.4f}")

# --- ANOVA a un facteur ---
# Rendement de 3 types d'engrais
engrais_A = [20.1, 21.5, 19.8, 22.0, 20.5, 21.2]
engrais_B = [23.4, 24.1, 22.8, 23.9, 24.5, 23.2]
engrais_C = [19.5, 20.2, 18.9, 20.8, 19.1, 20.5]

f_stat, p_val = stats.f_oneway(engrais_A, engrais_B, engrais_C)
print(f"\nANOVA: F={f_stat:.4f}, p={p_val:.6f}")

# Avec statsmodels (plus de details)
df_anova = pd.DataFrame({
    'rendement': engrais_A + engrais_B + engrais_C,
    'engrais': ['A']*6 + ['B']*6 + ['C']*6
})
model_anova = ols('rendement ~ C(engrais)', data=df_anova).fit()
anova_table = sm.stats.anova_lm(model_anova, typ=2)
print("\nTableau ANOVA:")
print(anova_table)

# Comparaisons multiples (Tukey)
from statsmodels.stats.multicomp import pairwise_tukeyhsd
tukey = pairwise_tukeyhsd(df_anova['rendement'], df_anova['engrais'], alpha=0.05)
print("\nComparaisons de Tukey:")
print(tukey)

# --- Diagnostic regression multiple ---
fig, axes = plt.subplots(1, 2, figsize=(12, 5))
residuals = model.resid
fitted = model.fittedvalues

axes[0].scatter(fitted, residuals, alpha=0.6)
axes[0].axhline(0, color='red', linestyle='--')
axes[0].set_xlabel('Valeurs ajustees')
axes[0].set_ylabel('Residus')
axes[0].set_title('Residus vs Ajustees')

stats.probplot(residuals, plot=axes[1])
axes[1].set_title('QQ-plot des residus')

plt.tight_layout()
plt.savefig("ch09_diagnostic.pdf")
plt.show()
\end{minted}
\end{python}

% ------------------------------------------------------------
\section{Exercices}
% ------------------------------------------------------------

\begin{exercice}[$\star$ -- R\'egression multiple]
Ajuster un mod\`ele de r\'egression du prix d'une voiture en fonction de l'ann\'ee, du kilom\'etrage et de la puissance (donn\'ees simul\'ees ou r\'eelles). Interpr\'eter les coefficients.
\end{exercice}

\begin{exercice}[$\star\star$ -- D\'emonstration]
Montrer que $\hat{\bm{\beta}}=(\mathbf{X}^\top\mathbf{X})^{-1}\mathbf{X}^\top\mathbf{Y}$ et que $\mathbf{H}=\mathbf{X}(\mathbf{X}^\top\mathbf{X})^{-1}\mathbf{X}^\top$ est idempotente et sym\'etrique.
\end{exercice}

\begin{exercice}[$\star\star$ -- ANOVA]
Quatre m\'ethodes p\'edagogiques sont test\'ees sur des groupes de 10 \'etudiants. R\'ealiser l'ANOVA, v\'erifier les hypoth\`eses (normalit\'e, homog\'en\'eit\'e) et effectuer les comparaisons de Tukey.
\end{exercice}

\begin{exercice}[$\star\star\star$ -- Projet : s\'election de mod\`ele]
Sur un jeu de donn\'ees r\'eel, comparer les mod\`eles par AIC, BIC et $R^2_{\mathrm{adj}}$. Utiliser la validation crois\'ee pour \'evaluer le pouvoir pr\'edictif.
\end{exercice}

% ------------------------------------------------------------
\begin{formules}
\begin{align*}
\hat{\bm{\beta}} &= (\mathbf{X}^\top\mathbf{X})^{-1}\mathbf{X}^\top\mathbf{Y} \\
\Cov(\hat{\bm{\beta}}) &= \sigma^2(\mathbf{X}^\top\mathbf{X})^{-1} \\
F_{\text{global}} &= \frac{R^2/(p-1)}{(1-R^2)/(n-p)} \\
R^2_{\mathrm{adj}} &= 1 - \frac{n-1}{n-p}(1-R^2) \\
F_{\text{ANOVA}} &= \frac{\mathrm{CME}}{\mathrm{CMR}} \sim F(k-1,n-k)
\end{align*}
\end{formules}
